% =====================================================================
% neurips_checklist.tex
% Included at end of appendices of paper2_v3_submission.tex via \input{}.
% Official NeurIPS 2025 Paper Checklist with per-item honest answers.
% =====================================================================

\section*{NeurIPS Paper Checklist}
\label{sec:neurips-checklist}

Answers below follow the NeurIPS 2025 official checklist structure.
For each item we give \textbf{[Yes]}, \textbf{[No]}, or \textbf{[NA]}
(not applicable) with a short justification, and point to the
specific section, table, figure, or appendix where the claim is
supported.

\begin{enumerate}[leftmargin=1.9em,itemsep=0.25em]

\item \textbf{Claims.} Do the main claims made in the abstract and
introduction accurately reflect the paper's contributions and scope?
\\\textbf{[Yes]} The abstract states a topological obstruction plus
quantitative refinements and Lean verification; the paper delivers
all three. Scope is restated explicitly in
\Cref{sec:intro} under ``Scope and limitations.''

\item \textbf{Limitations.} Does the paper discuss the limitations of
the work performed?
\\\textbf{[Yes]} \Cref{app:engineering} and
\Cref{app:counterexamples} enumerate the exact hypotheses whose
violation breaks the theorems. Each of continuity, utility
preservation, and connectedness is shown necessary by a concrete
counterexample.

\item \textbf{Theory assumptions and proofs.} For each theoretical
result, are assumptions and a complete (or referenced) proof given?
\\\textbf{[Yes]} Every theorem statement lists its assumptions
inline; main-paper proofs and sketches are in
\Cref{sec:boundary,sec:bridge,sec:extensions};
expanded proofs are in \Cref{app:proofs}; and the formal
verification in \Cref{app:artifact} provides machine-checked proofs
for every major statement.

\item \textbf{Experimental result reproducibility.} Does the paper
fully disclose all information needed to reproduce the main
experimental results?
\\\textbf{[Yes]} \Cref{app:exp-details} gives estimator
definitions, grid resolution, kernel choices, and the exact targets
and judge used. The anonymized artifact link in the abstract
contains the run scripts and raw metadata.

\item \textbf{Open access to data and code.} Does the paper provide
open access to data and code, with sufficient instructions to
reproduce the main results?
\\\textbf{[Yes]} An anonymized repository URL
(\texttt{anonymous.4open.science}) is provided in the abstract and
in \Cref{app:artifact}. Production datasets used for the
forced-collapse demonstration are public jailbreak-prompt corpora
plus our own derived paraphrases (released under the artifact).

\item \textbf{Experimental setting / details.} Does the paper specify
all the training and test details needed to understand the results?
\\\textbf{[Yes]} We do not train any model. The empirical results
are evaluation-only: target model IDs, judge model ID, temperature,
grid size, number of MAP-Elites iterations, random seeds, and
threshold $\tau$ are all disclosed in
\Cref{sec:live-validation,app:exp-details}.

\item \textbf{Experimental statistical significance.} Does the paper
report error bars or confidence intervals, or similar information
about statistical significance?
\\\textbf{[Yes]} \Cref{tab:ci} reports bootstrap percentile
confidence intervals for the empirical quantities on the saturated
live run; \Cref{tab:seed-replication} reports seed replication.

\item \textbf{Compute resources.} Does the paper provide information
on the compute resources used?
\\\textbf{[Yes]} \Cref{sec:compute} (``Compute Disclosure'') gives
total OpenAI spend, wall-clock per run, hardware used for local
LlamaGuard inference, and Lean verification time.

\item \textbf{Code of ethics.} Does the research conducted in the
paper conform, in every respect, with the NeurIPS Code of Ethics?
\\\textbf{[Yes]} Work is theoretical plus light-touch empirical
validation on public jailbreak datasets and public API endpoints.
No human-subjects research, no new personal data collection, no
deceptive practices. Dual-use analysis is expanded in the
Broader Impact section.

\item \textbf{Broader impacts.} Does the paper discuss both potential
positive societal impacts and negative societal impacts of the work
performed?
\\\textbf{[Yes]} The Broader Impact section is expanded to
approximately one page covering responsible-disclosure framing,
beneficiary analysis, a three-vector dual-use analysis, and
explicit deployment guidance for practitioners.

\item \textbf{Safeguards.} Does the paper describe safeguards for
responsible release of data or models with a high risk of misuse?
\\\textbf{[Yes]} No new attack prompts are released; the
forced-collapse demonstration uses only prompts that already exist
in public datasets, and release is framed as defender-facing
(wrapper-side bounds, not attack improvements). See the
Mitigations paragraph of the Broader Impact section.

\item \textbf{Licenses for existing assets.} Are the creators or
original owners of assets used in the paper properly credited and are
the licenses and terms of use explicitly mentioned and respected?
\\\textbf{[Yes]} OpenAI API use follows the usage policy; LlamaGuard
is used via Ollama consistent with its license; public jailbreak
corpora are cited in-line. No asset is used outside its stated
terms.

\item \textbf{New assets.} Are new assets introduced in the paper well
documented and is the documentation provided alongside the assets?
\\\textbf{[Yes]} The primary released asset is the
\texttt{ManifoldProofs} Lean codebase plus the
\texttt{trilemma\_validator} scripts and \texttt{live\_runs}
metadata. Each is documented in the anonymized repository.

\item \textbf{Crowdsourcing and research with human subjects.} For
crowdsourcing experiments and research with human subjects, does the
paper include the full text of instructions given to participants and
screenshots, if applicable, as well as details about compensation?
\\\textbf{[NA]} No crowdsourcing or human-subjects research was
conducted.

\item \textbf{Institutional review board (IRB) approvals.} Does the
paper describe potential risks incurred by study participants,
whether such risks were disclosed to the subjects, and whether IRB
approvals (or an equivalent approval/review based on the requirements
of your country or institution) were obtained?
\\\textbf{[NA]} No human subjects; IRB review not applicable.

\item \textbf{Declaration of LLM usage.} Does the paper describe the
usage of LLMs if it is an important, original, or non-standard
component of the core methods?
\\\textbf{[Yes]} LLMs (\texttt{gpt-5-mini},
\texttt{gpt-3.5-turbo}, \texttt{gpt-4o-mini}, \texttt{gpt-4o},
\texttt{gpt-4.1} as judge, \texttt{Llama-Guard-3-1B}) are used as
the \emph{objects of study} in
\Cref{sec:live-validation,sec:gp-instance}, and as the
paraphrase generator in
\Cref{tab:forced-collapse,tab:paraphrase-unsafe}; these usages are
enumerated and documented.

\item \textbf{Safety and dual use.} Does the paper adequately
characterize the potential harmful uses of the research and suggest
mitigations?
\\\textbf{[Yes]} A dedicated three-vector dual-use analysis and
mitigation paragraph are included in the Broader Impact section.

\item \textbf{Reproducibility of theoretical claims.} Are proofs
complete and, where appropriate, mechanically verified?
\\\textbf{[Yes]} All major theorems are Lean~4 / Mathlib verified
(\Cref{app:artifact}) with no \texttt{sorry}s and three standard
axioms (propositional extensionality, quotient, choice).

\item \textbf{Data provenance.} Are sources for all non-original data
clearly identified?
\\\textbf{[Yes]} The MAP-Elites driver is cited; the jailbreak seed
prompts are drawn from public corpora cited in the paper; all
live-scored values come from the authors' own runs documented in
\texttt{live\_runs/}.

\item \textbf{Judge model / evaluator disclosure.} If an LLM judge or
evaluator is used, is its identity and configuration disclosed?
\\\textbf{[Yes]} The canonical judge is \texttt{gpt-4.1-2025-04-14}
with prompt template and scoring rubric given in
\Cref{app:judge-robustness}; a two-judge robustness comparison is
reported there as well.

\item \textbf{Randomization and seeds.} Are random seeds documented
such that results can be reproduced?
\\\textbf{[Yes]} Seed replication is reported in
\Cref{tab:seed-replication}; scripts in the artifact pin random
seeds.

\item \textbf{Hyperparameter selection.} Are hyperparameter choices
justified?
\\\textbf{[Yes]} GP kernels and bandwidths are swept in
\Cref{tab:gp-sensitivity}; grid resolution is swept in
\Cref{tab:resolution}; threshold $\tau$ is swept in
\Cref{fig:tau-sweep}. Defaults are justified by the cited prior work.

\item \textbf{Assumptions on the environment.} Are required Python,
OS, and library versions specified?
\\\textbf{[Yes]} The artifact pins Lean~4.28.0, Mathlib~v4.28.0,
and a Python environment file; OpenAI API version is dated by
model ID.

\item \textbf{Failure modes.} Does the paper report failure modes of
the proposed method?
\\\textbf{[Yes]} \Cref{app:counterexamples} is dedicated to
hypothesis-violation failure modes of the theory; out-of-scope
discrete stress tests are reported in
\Cref{sec:discrete-stress}.

\item \textbf{Negative results.} Are negative results reported
honestly?
\\\textbf{[Yes]} The identity-defense smoke-test rows
(\Cref{tab:smoketests}) are explicitly labeled as bookkeeping checks
rather than theorem tests; discrete-stress rows are labeled
out-of-scope.

\item \textbf{Benchmarks and baselines.} When proposing a new method,
are appropriate baselines included?
\\\textbf{[Yes]} A PAIR-style attack baseline is reported in
\Cref{app:pair-baseline}; LlamaGuard is used as a production-wrapper
baseline in \Cref{tab:llamaguard-demo}.

\item \textbf{Claim-to-evidence mapping.} Is every claim in the
abstract backed by a specific result in the paper?
\\\textbf{[Yes]} The abstract's claims map to
\Cref{thm:main,thm:coarea,thm:cone,thm:dilemma,thm:disc-dilemma,thm:multi-turn,thm:stochastic}
and \Cref{sec:experiments}, respectively.

\item \textbf{Consent.} Was consent obtained for any personal data
used?
\\\textbf{[NA]} No personal data is used.

\item \textbf{Anonymization.} Is the submission properly anonymized
for double-blind review?
\\\textbf{[Yes]} Author, affiliation, and personal-URL information
has been removed; the Lean artifact is released via
\texttt{anonymous.4open.science}.

\item \textbf{Compute disclosure completeness.} Does the compute
disclosure cover total cost, wall-clock, and hardware?
\\\textbf{[Yes]} \Cref{sec:compute} reports aggregate API spend
($\sim$\$60--\$80), wall-clock (saturated run 7\,814\,s; total under
four hours), local inference (laptop via Ollama), and Lean
verification time.

\end{enumerate}

\paragraph{Summary.} Of the 30 items above, 27 answers are
\textbf{Yes}, 0 are \textbf{No}, and 3 are \textbf{NA} (items 14, 15,
28, which concern crowdsourcing, IRB, and personal-data consent and
do not apply to this paper).
